\documentclass[11pt,a4paper]{article}
\usepackage[margin=1in]{geometry}
\usepackage{amsmath,amssymb,amsthm}
\usepackage{algorithm}
\usepackage{algpseudocode}
\usepackage{booktabs}
\usepackage{hyperref}

\theoremstyle{definition}
\newtheorem{definition}{Definition}[section]
\newtheorem{theorem}{Theorem}[section]
\newtheorem{lemma}[theorem]{Lemma}
\newtheorem{remark}{Remark}[section]

\title{\textbf{Rigorous Mathematical Specification, Gradient Highway Dynamics, and Signal Propagation of Residual Skip Connections}}
\author{Team Scaibu \\ \small Email: \texttt{jaydeep.9664920749@gmail.com} \\ \small Architecture and Core Engine Team}
\date{October 2026}

\begin{document}
\maketitle

\begin{abstract}
This document presents the complete mathematical formulation and algorithmic specification of Residual Skip Connections ($y = h(x) + \alpha F(x)$). By creating an additive identity highway through deep computation graphs, residual networks eliminate exponential gradient decay ($\prod_{l=1}^L W_l \to 0$), enabling the stable optimization of arbitrarily deep networks ($L > 1000$). We formalize the gradient unrolling identity, formulate the complete algorithmic pseudo-code matching the reference implementation, derive exact computational complexities, step through a numerical example, and analyze dimension projection and gating.
\end{abstract}

\section{Notation and Mathematical Preliminaries}

Let $x \in \mathbb{R}^{B \times D_{\text{in}}}$ denote an input identity tensor, and let $\mathcal{F}(x) \in \mathbb{R}^{B \times D_{\text{out}}}$ denote the non-linear sub-layer transformation.

\begin{itemize}
    \item $W_{\text{proj}} \in \mathbb{R}^{D_{\text{out}} \times D_{\text{in}}}$: Optional projection shortcut matrix for dimension matching.
    \item $\alpha > 0$: Branch residual scaling coefficient (default $1.0$).
    \item $g \in \mathbb{R}^{D_{\text{out}}}$: Optional channel-wise gating parameter vector.
    \item $y \in \mathbb{R}^{B \times D_{\text{out}}}$: Output merged activation tensor.
\end{itemize}

\begin{table}[h]
\centering
\caption{Summary of Mathematical Symbols for Residual Connection}
\begin{tabular}{lll}
\toprule
\textbf{Symbol} & \textbf{Type / Domain} & \textbf{Description} \\
\midrule
$x$ & $\mathbb{R}^{B \times D_{\text{in}}}$ & Input skip tensor \\
$\mathcal{F}(x)$ & $\mathbb{R}^{B \times D_{\text{out}}}$ & Sub-layer residual branch output \\
$W_{\text{proj}}$ & $\mathbb{R}^{D_{\text{out}} \times D_{\text{in}}}$ & Shortcut projection matrix \\
$\alpha$ & $\mathbb{R}_{> 0}$ & Residual scaling factor \\
$y$ & $\mathbb{R}^{B \times D_{\text{out}}}$ & Combined output tensor \\
\bottomrule
\end{tabular}
\end{table}

\section{Problem Definition and Core Concepts}

\subsection{Intuition and Motivation}
In plain deep networks ($x_{l+1} = \mathcal{F}_l(x_l)$), the gradient backpropagating from layer $L$ to layer $l$ is governed by the chain rule product $\frac{\partial \mathcal{L}}{\partial x_l} = \frac{\partial \mathcal{L}}{\partial x_L} \prod_{k=l}^{L-1} \frac{\partial \mathcal{F}_k}{\partial x_k}$. As $L \to \infty$, this product vanishes or explodes exponentially. In residual networks ($x_{l+1} = x_l + \mathcal{F}_l(x_l)$), unrolling gives $x_L = x_l + \sum_{k=l}^{L-1} \mathcal{F}_k(x_k)$, which yields $\frac{\partial \mathcal{L}}{\partial x_l} = \frac{\partial \mathcal{L}}{\partial x_L} \left( I + \frac{\partial}{\partial x_l} \sum_{k=l}^{L-1} \mathcal{F}_k(x_k) \right)$. The additive identity term $I$ ensures that gradients can flow directly back to initial layers without vanishing.

\subsection{Formal Mathematical Definition}
\begin{definition}[Residual Topology Equations]
\begin{itemize}
    \item \textbf{Standard Residual}: $y = x + \mathcal{F}(x)$
    \item \textbf{Scaled Residual}: $y = x + \alpha \mathcal{F}(x)$
    \item \textbf{Projection Residual}: $y = x W_{\text{proj}}^T + \alpha \mathcal{F}(x)$
    \item \textbf{Gated Residual}: $y = (1 - \sigma(g)) \odot x + \sigma(g) \odot \mathcal{F}(x)$
\end{itemize}
\end{definition}

\section{Algorithmic Specification}

The computational pipeline implemented in \texttt{NnAlgoResidualConnection} is detailed below.

\begin{algorithm}[H]
\caption{Residual Connection Forward Merge}
\begin{algorithmic}[1]
\Require Identity batch $X \in \mathbb{R}^{B \times D_{\text{in}}}$, sublayer batch $F \in \mathbb{R}^{B \times D_{\text{out}}}$
\Require Optional projection $W_{\text{proj}} \in \mathbb{R}^{D_{\text{out}} \times D_{\text{in}}}$, scale $\alpha > 0$, topology $\in \{\text{standard\_add}, \text{scaled\_add}, \text{gated\_residual}\}$
\Ensure Combined output tensor $Y \in \mathbb{R}^{B \times D_{\text{out}}}$
\State $B \gets \text{rows}(X), \quad D_{\text{in}} \gets \text{cols}(X), \quad D_{\text{out}} \gets \text{cols}(F)$
\State $Y \gets \text{new Matrix of shape } (B, D_{\text{out}})$
\State $H \gets \text{new Matrix of shape } (B, D_{\text{out}})$
\If{$W_{\text{proj}} \text{ is not null}$}
    \State $H \gets X \cdot W_{\text{proj}}^T$
\ElsIf{$D_{\text{in}} = D_{\text{out}}$}
    \State $H \gets X$
\Else
    \State \textbf{throw} PreconditionFailureException(``Dimension mismatch without projection matrix'')
\EndIf
\For{$b \gets 1 \textbf{ to } B$}
    \For{$d \gets 1 \textbf{ to } D_{\text{out}}$}
        \If{$\text{topology} = \text{standard\_add}$}
            \State $Y[b][d] \gets H[b][d] + F[b][d]$
        \ElsIf{$\text{topology} = \text{scaled\_add}$}
            \State $Y[b][d] \gets H[b][d] + \alpha \cdot F[b][d]$
        \EndIf
    \EndFor
\EndFor
\State \Return $\{\text{output\_tensor}: Y\}$
\end{algorithmic}
\end{algorithm}

\section{Theoretical Guarantees and Gradient Highway Preservation}

\begin{theorem}[Additive Gradient Highway Non-Vanishing Property]
For any deep residual network $x_L = x_l + \sum_{k=l}^{L-1} \mathcal{F}_k(x_k)$, the backward gradient satisfies:
\begin{equation}
\frac{\partial \mathcal{L}}{\partial x_l} = \frac{\partial \mathcal{L}}{\partial x_L} + \sum_{k=l}^{L-1} \frac{\partial \mathcal{L}}{\partial x_L} \frac{\partial \mathcal{F}_k}{\partial x_l}
\end{equation}
The gradient $\frac{\partial \mathcal{L}}{\partial x_l}$ cannot vanish uniformly for all layers $l$, even when individual weight matrices $\frac{\partial \mathcal{F}_k}{\partial x_k}$ are small.
\end{theorem}

\subsection{Limitations}
\begin{itemize}
    \item \textbf{Residual Stream Variance Expansion}: Without scaling $\alpha = 1/\sqrt{2L}$, activation variance grows as $\mathcal{O}(L)$ (addressed by ALGO-NN-49).
\end{itemize}

\section{Complexity Analysis}

\subsection{Time and Space Complexity}
\begin{itemize}
    \item \textbf{Time}: Without projection, $\Theta(B \cdot D_{\text{out}})$ scalar additions. With projection, $\Theta(B \cdot D_{\text{in}} \cdot D_{\text{out}})$ FLOPs.
    \item \textbf{Space}: $\Theta(B \cdot D_{\text{out}})$ output tensor memory.
\end{itemize}

\section{Worked Numerical Example}

Consider $B=1, D=2$, $x = [1.0, 2.0]$, $F = [0.5, -0.5]$, $\alpha = 0.5$, $\text{topology} = \text{scaled\_add}$.

\begin{align}
y_1 &= 1.0 + 0.5(0.5) = 1.0 + 0.25 = 1.25 \\
y_2 &= 2.0 + 0.5(-0.5) = 2.0 - 0.25 = 1.75
\end{align}
Output: $Y = [1.25, 1.75]$.

\section{Numerical Considerations and Edge Cases}

\begin{itemize}
    \item \textbf{In-place Addition}: When $D_{\text{in}} = D_{\text{out}}$, executing $X \leftarrow X + \alpha F$ in-place minimizes GPU memory allocations during forward/backward passes.
\end{itemize}

\begin{thebibliography}{9}
\bibitem{he2016deep} K.~He, X.~Zhang, S.~Ren, and J.~Sun, ``Deep residual learning for image recognition,'' in \emph{IEEE Conference on Computer Vision and Pattern Recognition (CVPR)}, pp.~770--778, 2016.
\bibitem{he2016identity} K.~He et al., ``Identity mappings in deep residual networks,'' in \emph{European Conference on Computer Vision (ECCV)}, pp.~630--645, 2016.
\end{thebibliography}

\end{document}
